\section{Inference}

Based on the inference algorithms in propositional logics we now derive generalized versions for CSPs.

%\subsection{Disentangling}
%



\subsection{Solution by disentanglement}

Let us now use the manipulation approaches to derive solution approaches to CSPs.

\subsubsection{Elementarization}

Constraint Satisfaction Problems on elementary constraint networks have solutions by cartesian products of the local supports.

We want to find elementary constraint networks which support is included in the support of the constraint network.

We notice that finding a single solution assignment to a constraint network $\extnet$ is finding a one-hot encoding with
\begin{align*}
    \onehotmapofat{\catindexof{\nodes}}{\nodevariables}
    \prec
    \contractionof{\extnet}{\nodevariables} \, .
\end{align*}

One-hot encodings are in particular non-vanishing elementary tensors.

Let us generalize this approach to solving CSPs by elementarization:
\begin{itemize}
    \item Find a non-vanishing elementary tensor network, which entails the
\end{itemize}

Modus ponens is an elementarization approach: The implication edges are split into elementary vectors.

\subsubsection{Vanishing sub-contraction}

When a sub-contraction of the constraint network is vanishing, we know that the CSP has no solution.

This is a typical approach in logics, e.g. the aim of resolution.

\subsection{Generlized unit propagation}

%The generalization of
%\subsubsection{Generalized Forward Chaining}

We generalize forward chaining to arbitrary constraint networks.
Main elements:
\begin{itemize}
    \item Constraint vectors are contracted with constraint tensors containing the corresponding node.
    \item If the split successes and results in a new constraint vector, we continue step 1 with the new constraint vector.
\end{itemize}

\begin{algorithm}[hbt!]
    \caption{Generalized Forward Chaining}\label{alg:genForwardChaining}
    \begin{algorithmic}
        \Require Constraint tensor network $\extnet$ on $\graph=(\nodes,\edges)$, message edges $\domainedges$ and a set $\arbset$ of subsets of $\edges$ for local propagation
        \Ensure
        \iosepline
        \State
        \State Initialize a queue $\graphqueue$ by all nodes with constraint vectors
%        \begin{align*}
%            \graphqueue =
%        \end{align*}
        \While{$\graphqueue$ is not empty}
            \State Choose a node variable $\node$ from $\graphqueue$
            \For{All $\edge\in\domainedges$ with $\node\in\edge$}
                \State Contract
                \begin{align*}
                    \hypercoreofat{\edge}{\catvariableof{\edge}}
                    = \contractionof{\hypercoreofat{\edge}{\catvariableof{\edge}},\hypercoreofat{\{\node\}}{\catvariableof{\node}}}{\catvariableof{\edge}}
                \end{align*}
                \State Try to split $\hypercoreofat{\edge}{\catvariableof{\edge}}$ into $\{\node\}$ and $\edge\backslash\{\node\}$
                \If{Splitting successful and $\cardof{\edge\backslash\{\node\}}=1$}
                    \State Add $\edge\backslash\{\node\}$ to the queue
                \EndIf
            \EndFor
        \EndWhile
        \State \Return $\{\hypercoreofat{\edge}{\catvariableof{\edge}}\wcols\edge\in\domainedges\}$
    \end{algorithmic}
\end{algorithm}

\subsection{Constraint Propagation}

\todo{
    Understand constraint propagation as an iterative application of the $(i)$ manipulation rule where messages are computed and contracted into memory cores.
    (Each step applies the rule twice: New message core computed and fused into old message core.)
    Call $\kcore$ memory core, which is updated by a message.
}

%\subsect{Message cores}

To store preliminary conclusions, we define auxiliary message cores storing constraints on variables $\edge\subset\nodes$.
They are understood as logical formulas to the atomization variables $\left(\catvariableof{\edge}=\catindexof{\edge}\right)$ of the respective formulas
\begin{align*}
    \kcoreofat{\edge}{\catvariableof{\edge}}
    = \bigvee_{\catindexof{\edge} \wcols \kcoreofat{\edge}{\indexedcatvariableof{\edge}}} \bigwedge_{\node\in\edge} \left(\catvariableof{\edge}=\catindexof{\edge}\right) \, .
\end{align*}

\begin{definition}\label{def:knowledgeCoreSoundComplete}
    Let $\extnet$ be a constraint satisfaction problem. % and $$ be a message core. % \left(\catvariableof{\edge}=\catindexof{\edge}\right)
    We say that a message core $\kcoreofat{\edge}{\catvariableof{\edge}}$ is sound for $\extnet$, if
    \begin{align*}
        \nonzerocirc\contractionof{\extnet}{\catvariableof{\edge}}  \prec \kcoreofat{\edge}{\catvariableof{\edge}}
    \end{align*}
    and complete for $\extnet$ if in addition
    \begin{align*}
        \nonzerocirc\contractionof{\extnet}{\catvariableof{\edge}}=\kcoreofat{\edge}{\catvariableof{\edge}} \, .
    \end{align*}
\end{definition}

We now provide a solution algorithm for constraint satisfaction problems by propagating local contractions.
The dynamic programming paradigm is implemented by the storage of partial entailment results in message cores.
We then iterate over local entailment checks, where we recursively add further entailment checks to be redone due to additional knowledge.
This local entailment scheme is called Knowledge Propagation and described in a generic way in \algoref{alg:knowledgePropagation}.

\red{
    We understand the support of local contractions as messages, which are tensors $\kcoreofat{\edge}{\catvariableof{\edge}}$ associated with message edges $\edge$.
    These message cores are propagated through the network, and summarize neighboring constraints.
}

\todo{
    Adjust constraint propagation to match tnreason message passing schemes:
    Queue containing $\secgraph$ and open nodes $\secnodes$.
}

\begin{algorithm}[hbt!]
    \caption{Constraint Propagation}\label{alg:knowledgePropagation}
    \begin{algorithmic}
        \Require Constraint tensor network $\extnet$ on $\graph=(\nodes,\edges)$, message edges $\domainedges$ and a set $\arbset$ of subsets of $\edges$ for local propagation
        \Ensure Message cores $\kcoreofat{\edge}{\catvariableof{\edge}}$ for $\edge\in\domainedges$ with $\contractionof{\extnet}{\catvariableof{\edge}}\prec\kcoreofat{\edge}{\catvariableof{\edge}}$
        \iosepline
        \State
        \State Initialize for all $\edge\in\domainedges$:
        \begin{align*}
            \kcoreofat{\edge}{\catvariableof{\edge}}
            =\onesat{\catvariableof{\edge}}
        \end{align*}
        \State Initialize a queue
        \begin{align*}
            \graphqueue =
            \left\{
                (\edgesof{inference},\edgeof{m}) \wcols \edgesof{inference}\in\arbset \ncond \edgeof{m}\in\domainedges \ncond \edgeof{m} \cap \big(\bigcup_{\edge\in\edgesof{inference}}\edge\big) \neq \varnothing
            \right\}
        \end{align*}
        \While{$\graphqueue$ is not empty}
            \State Choose a set of edges from the queue
            \begin{align*}
                \inferenceedges,\edgeof{m} \algdefsymbol \graphqueue\mathrm{.pop()}
            \end{align*}
            %\ForAll{$\edge\in\domainedges$ with $\edge\cap\bigcup_{\secedge\in\secedges}\secedge\neq\varnothing$}
            \State Contract
            \begin{align*}
                \hypercoreat{\catvariableof{\messageedge}}
                = \nonzerocirc
                \contractionof{
                    \{\hypercoreofat{\secedge}{\catvariableof{\secedge}}\wcols\secedge\in\inferenceedges\}
                    \cup \{\kcoreofat{\edge}{\catvariableof{\edge}}\wcols\edge \in \domainedges , \, \edge\cap\bigcup_{\secedge\in\secedges}\secedge\neq\varnothing \}
                }{\catvariableof{\messageedge}}
            \end{align*}

%            \begin{align*}
%                \hypercoreat{\catvariableof{\edge}}
%                = \nonzerocirc\contractionof{
%                    \{\hypercoreofat{\secedge}{\catvariableof{\secedge}}\wcols\secedge\in\secedges\}
%                    \cup \{\kcoreofat{\edge}{\catvariableof{\edge}}\wcols\edge \in \domainedges , \, \edge\cap\bigcup_{\secedge\in\secedges}\secedge\neq\varnothing \}
%                }{\catvariableof{\edge}}
%            \end{align*}
            \If{$\hypercoreat{\catvariableof{\edge}}\neq\kcoreofat{\edge}{\catvariableof{\edge}}$}
                \begin{align*}
                    \kcoreofat{\edge}{\catvariableof{\edge}} \algdefsymbol \hypercoreat{\catvariableof{\edge}}
                \end{align*}
                \ForAll{$\secedges\in\arbset$ with $\edge\cap\bigcup_{\secedge\in\secedges}\secedge\neq\varnothing$}
                    \begin{align*}
                        \graphqueue\mathrm{.push(}\secedges\mathrm{)}
                    \end{align*}
                \EndFor
            \EndIf
        %\EndFor
        \EndWhile
        \State \Return $\{\kcoreofat{\edge}{\catvariableof{\edge}}\,:\,\edge\in\domainedges\}$
    \end{algorithmic}
\end{algorithm}

% Interpretation
Each chosen subset $\secedges\in\arbset$ is understood as a local knowledge base, which is then applied for local entailment.
The message cores are understood as messages, which propagate information from different regions of a tensor network (see \charef{cha:messagePassing}).

%Implementation
There are different ways of implementing \algoref{alg:knowledgePropagation}, by choosing the set $\arbset$ of constraint sets $\secedges$ and domain $\domainedges$. % and a stopping criterion.


\begin{example}[AC-3 as a Constraint propagation scheme]
    The AC-3 algorithm (see \cite{mackworth_consistency_1977}) is a specific instance, where message cores are assigned to single variables and propagation is performed on single constraint cores.
    We assume that each edge has cardinality $2$, that is $\graph$ is a graph.
    Here the inference clusters are
    \begin{align*}
        \arbset
        = \big\{\{\edge\} \wcols \edgein\big\} \, .
    \end{align*}
    The message edges are
    \begin{align*}
        \domainedges
        = \big\{\{\node\} \wcols \nodein\big\} \, .
    \end{align*}
    At each local inference step to the edge $\edge=\{\nodeof{0},\nodeof{1}\}$ we compute
    \begin{align*}
        \contractionof{
            \kcoreofat{\{\nodeof{0}\}}{\catvariableof{\nodeof{0}}},
            \hypercoreofat{\edge}{\catvariableof{\nodeof{0},\nodeof{1}}},
            \kcoreofat{\{\nodeof{0}\}}{\catvariableof{\nodeof{1}}}
        }{\catvariableof{\nodeof{0}}} \, .
    \end{align*}


\end{example}


%The central property used in constraint propagation is that any subcontraction can be added to the constraint network without changing it.
%\begin{theorem}\label{the:booleanContractionInvariance}
%	Given a boolean tensor network $\extnet$ on $\graph$ and $\secedges\subset\edges$.
%	For $\secgraph=(\secnodes,\secedges)$ with $\nodes=\cup_{\edge\in\secedges}\edge$ and the tensor network $\tnetof{\secgraph}$ with boolean tensors coinciding on $\secedges$ with those in $\extnet$ we have
%	\begin{align*}
%		\contractionof{\extnet}{\catvariableof{\nodes}} =
%		\contractionof{\extnet\cup\{\nonzerocirc\contractionof{\tnetof{\secgraph}}{\thirdnodes}\}}{\catvariableof{\nodes}} \, ,
%	\end{align*}
%	where $\thirdnodes\subset\nodes$ is arbitrary.
%\end{theorem}
%\begin{proof}
%	We will proof this statement later in \charef{cha:messagePassing}, see \theref{the:invarianceAddingSubcontractions}.
%\end{proof}

\begin{theorem}
    \label{the:soundnessKnowledgePropagation}
    At any state of the constraint propagation \algoref{alg:knowledgePropagation}, we have that each message core $\kcoreof{\secedge}$ is sound for $\extnet$.
    After each update in \algoref{alg:knowledgePropagation}, $\kcoreof{\secedge}$ is further monotonically decreasing with respect to the partial ordering.
\end{theorem}
\begin{proof}
    \todo{This follows from soundness of the $i$ manipulation rule.}
    We show the first claim by induction over the update steps in \algoref{alg:knowledgePropagation}.
    At the start, where $\kcoreofat{\edge}{\edgevariables} = \onesat{\edgevariables}$, we trivially have
    \begin{align*}
        \contractionof{\extnet \cup \{\kcoreofat{\edge}{\edgevariables} \wcols \edge\in\domainedges\}}{\nodevariables}
        = \contractionof{\extnet \cup \{\onesat{\edgevariables} \wcols \edge\in\domainedges\}}{\nodevariables}
        = \contractionof{\extnet}{\nodevariables} \, .
    \end{align*}
    Let us now assume, that for a state of cores $\{\kcoreof{\edge} \wcols \edge\in\domainedges\}$ the first claim holds and let $\secedges\subset\edges$ be chosen for the update of $\kcoreof{\secedge}$.
    By the invariance under adding the support of subcontractions, which we will proof in more detail as \theref{the:monotonicityBinaryContractions} in \charef{cha:messagePassing}, we have for the update
    \begin{align*}
        \tilde{\kcore}^{\secedge}[\catvariableof{\secedge}]
        = \nonzerocirc\contractionof{
            \{\hypercoreof{\edge} \wcols \edge\in\secedges\} \cup \{\kcoreof{\edge} \wcols \edge\in\domainedges, \, \edge\cap\bigcup_{\secedge\in\secedges}\secedge\neq\varnothing  \}
        }{\catvariableof{\secedge}}
    \end{align*}
    that
    \begin{align*}
        \contractionof{\extnet \cup \{\kcoreofat{\edge}{\edgevariables} \wcols \edge\in\domainedges\}}{\nodevariables}
        = \contractionof{\extnet \cup \{\kcoreofat{\edge}{\edgevariables} \wcols \edge\in\domainedges\} \cup \{\tilde{\kcore}^{\secedge}[\catvariableof{\secedge}]\}
        }{\nodevariables}  \, .
    \end{align*}
    Thus, the first claim holds also after the update of the core to $\secedge$.

    % Since the kcore is in the contraction
    We further have with the monotonicity of boolean contraction (see \theref{the:monotonicityBinaryContractions}) that for any update of $\kcoreof{\secedge}$ by $\tilde{\kcore}^{\secedge}$
    \begin{align*}
        \tilde{\kcore}^{\secedge}[\catvariableof{\secedge}]
        = \nonzerocirc\contractionof{
            \{\hypercoreof{\edge} \wcols \edge\in\secedges\}
            \cup \{\kcoreof{\edge} \wcols \edge\in\domainedges, \, \edge\cap\bigcup_{\edge\in\secedges}\edge\neq\varnothing\}
        }{\catvariableof{\secedge}}
        \prec \kcoreofat{\secedge}{\catvariableof{\secedge}} \, .
    \end{align*}
    Thus, each message core is monotonously decreasing at each update, with respect to the partial tensor ordering.

    From the first claim we further have for any $\secedge\in\secedges$
    \begin{align*}
        \contractionof{\{\kcoreofat{\secedge}{\catvariableof{\secedge}}\}\cup\extnet \cup \left\{\kcoreof{\edge} \wcols \edge\in\domainedges/\{\secedge\}\right\}}{\catvariableof{\secedge}}
        = \contractionof{\extnet}{\catvariableof{\secedge}}
    \end{align*}
    And thus in combination with the monotonicity of boolean contraction (see \theref{the:monotonicityBinaryContractions}) that
    \begin{align*}
        \nonzeroof{\contractionof{\extnet}{\catvariableof{\secedge}}} \prec \kcoreofat{\secedge}{\catvariableof{\secedge}}  \, . & \qedhere
    \end{align*}
\end{proof}

Let us now show that constraint propagation always terminates.
We can further characterize the message cores at termination.

\begin{definition}
    We say that a set of message cores $\{\kcoreof{\edge} \wcols \edge\in\domainedges\}$ is consistent with a set $\{\hypercoreof{\edge} \wcols \edge\in\secedges\}$, if for any $\edge\in\domainedges$
    \begin{align*}
        \kcoreofat{\edge}{\catvariableof{\edge}}
        = \nzcontractionof{\{\hypercoreof{\edge} \wcols \edge\in\secedges\} \cup \{\kcoreof{\edge} \wcols \edge\in\domainedges\}}{\catvariableof{\edge}} \, .
    \end{align*}
\end{definition}

This property is similar to the completeness of a message core, when interpreting the other message cores and the constraints $\{\hypercoreof{\edge} \wcols \edge\in\secedges\}$ as posing a Constraint Satisfaction Problem.

\begin{theorem}
    The Constraint propagation \algoref{alg:knowledgePropagation} always terminates.
%    At termination, each message core $\kcoreof{}$
    At termination we further for each $\secedges\in\arbset$ and $\edge\in\domainedges$ with $\edge\cap\bigcup_{\edge\in\secedges}\edge\neq\varnothing$, that the message cores $\{\kcoreof{\edge} \wcols \edge\in\domainedges, \, \edge\cap\bigcup_{\edge\in\secedges}\edge\neq\varnothing\}$ are consistent with $\{\hypercoreof{\edge} \wcols \edge\in\secedges\}$.
\end{theorem}
\begin{proof}
    For each message core, there are finitely many boolean tensor precessing it with respect to the partial order.
    Therefore, since they are monotonously decreasing, each message core can only be varied finitely many times during the algorithm.
    In total the algorithm can run only finitely many times in the second for loop, where new sets of edges are pushed into the queue.
    Therefore the while loop will always terminate.

    When after a single pass through the while loop with chosen $\secedges\in\arbset$, the set $\secedges$ is not pushed back into $\graphqueue$, we have for any $\edge\in\domainedges$ with $\edge\cap\bigcup_{\edge\in\secedges}\edge\neq\varnothing$ that
%    At termination, we have for any $\secedges\in\arbset$
    \begin{align*}
        \kcoreofat{\edge}{\catvariableof{\edge}}
        = \nzcontractionof{\{\hypercoreof{\edge} \wcols \edge\in\secedges\} \cup \{\kcoreof{\edge} \wcols \edge\in\domainedges, \, \edge\cap\bigcup_{\edge\in\secedges}\edge\neq\varnothing\}}{\catvariableof{\edge}} \, .
    \end{align*}
    Whenever the contraction on the right hand side changes during the algorithm, the set $\secedges$ is pushed into $\graphqueue$.
    At termination of the algorithm, $\graphqueue$ is empty, and the claimed consistency therefore has to hold.
%    since otherwise the queue $\graphqueue$ would at least contain $\secedges$ and the while loop would not terminate.
\end{proof}

%A direct consequence of this theorem is that when running constraint propagation with all edges, then
%\begin{corollary}
%    When we choose $\arbset=\{\edges\}$ for an CSP $\extnet$, then the message cores returned by \algoref{alg:knowledgePropagation} are complete.
%\end{corollary}


We can exploit the constraint propagation \algoref{alg:knowledgePropagation} for the solution of Constraint Satisfaction Problems, by taking $\extnet$ as the tensor network of constraint tensors.
Whenever a message core vanishes, we can conclude that the Constraint Satisfaction Problems is not satisfiable, as we show next.

\begin{corollary}
    Let us for a Constraint Satisfaction Problem encoded by $\extnet$ run constraint propagation \algoref{alg:knowledgePropagation}.
    Whenever for any $\secedge\in\edges$ we have $\kcoreofat{\secedge}{\catvariableof{\secedge}}=\zerosat{\catvariableof{\secedge}}$, then the Constraint Satisfaction Problem is not satisfiable.
\end{corollary}
\begin{proof}
    Whenever $\kcoreofat{\secedge}{\catvariableof{\secedge}}=\zerosat{\catvariableof{\secedge}}$, then we have by \theref{the:soundnessKnowledgePropagation}
    \begin{align*}
        \nonzeroof{\contractionof{\extnet}{\catvariableof{\secedge}}} \prec \zerosat{\catvariableof{\secedge}}
    \end{align*}
    and therefore
    \begin{align*}
        \contraction{\extnet} = 0 \, . & \qedhere
    \end{align*}
\end{proof}


\subsect{Applications}

Let us exemplify the usage of constraint propagation on Constraint Satisfaction Problems posed by entailment queries on Markov Networks.

\begin{corollary}
    \label{cor:knowledgePropagationMarkovNetworks}
    Let \algoref{alg:knowledgePropagation} be run on the cores $\tnetof{\graph}\cup\{\bencodingof{\exformula}\}$ with an arbitrary design of $\secedges$.
    Whenever for a formula $\formulaat{\catvariableof{\secnodes}}$ and a $\kcoreof{\edge}$ we have
    \[ \contractionof{\kcoreof{\edge},\bencodingof{\exformula}}{\exformulavar=0} =0  \]
    then the Markov Network $\extnet$ probabilistically entails $\exformula$.
    If on the contrary
    \[ \contractionof{\kcoreof{\edge},\bencodingof{\exformula}}{\exformulavar=1} =0  \]
    then the Markov Network $\extnet$ probabilistically entails $\lnot\exformula$, that is probabilistically contradicts $\exformula$.
\end{corollary}
\begin{proof}
    This follows from \theref{the:soundnessKnowledgePropagation} ensuring the soundness of constraint propagation and the sufficiency of local entailment.
\end{proof}

\begin{example}[Batch decision of entailment]
    Let $\formulaset$ be a set of formulas and $\probofat{\graph}{\shortcatvariables}$ a Markov Network, for which it shall be decided, which formulas in $\formulaset$ are entailed, contradicted or contingent.
    We can in addition to the cores of the Markov Network create the cores $\{\bencodingofat{\exformula}{\formulavar,\shortcatvariables} \wcols \exformula\in\formulaset\}$ and prepare the message cores
    \begin{align*}
        \kcoreofat{\{\formula\}}{\formulavar} \, .
    \end{align*}
    To decide entailment batchwise, constraint propagation \algoref{alg:knowledgePropagation} can be run.
    Whenever during the algorithm we have that for a $\formula$, then \corref{cor:knowledgePropagationMarkovNetworks} implies that if
    \begin{align*}
        \kcoreofat{\{\formula\}}{\formulavar} =
        \begin{cases}
            \tbasisat{\formulavar} & \text{then, the formula is entailed by }\probof{\graph} \, . \\
            \fbasisat{\formulavar} & \text{then, the formula is contradicted by }\probof{\graph} \, .\\
            \onesat{\formulavar} & \text{then no conclusion can be drawn.}
        \end{cases}
    \end{align*}
    Note, that $\kcoreofat{\{\formula\}}{\formulavar} = \zerosat{\formulavar}$ can not happen, since this would mean that $\nonzeroof{\probof{\graph}}$ is inconsistent.
    Thus, at any stage of \algoref{alg:knowledgePropagation}, one of the three holds.
    % Inference rules such as modus ponens can be mimicked.
\end{example}

\subsection{Backtracking search}

\todo{
    Understand as adding basis vectors which are locally possible (i.e. can satisfy affected constraints).
}

% Combination of constraint propagation with backtracking search
When the constraint propagation \algoref{alg:knowledgePropagation} converges in a given implementation and no message core vanishes, we can however not conclude that the Constraint Satisfaction Problem is not satisfiable.
However, for any index tuple $\catindexof{\nodes}$ to be a solution of the CSP to $\extnet$, we have the necessary condition
\begin{align*}
    \uniquantwrtof{\secedge\in\secedges}{\kcoreofat{\edge}{\catvariableof{\secedge}=\restrictionofto{\catindexof{\nodes}}{\secedges}}=1} \, ,
\end{align*}
where by $\restrictionofto{\catindexof{\nodes}}{\secedges}$ we denote the restriction of the index tuple $\catindexof{\nodes}$ to the variables included in $\secedges$.
One can use this insight as a starting point for backtracking search, where the assignments to variables $\catvariableof{\secnodes}$ are iteratively guessed, based on the restriction that each constraint is locally satisfialbe, i.e. .
\begin{align*}
    \uniquantwrtof{\secedge\in\secedges}{
        \contraction{\kcoreofat{\edge}{
            \catvariableof{\secedge\cap\secnodes}=\restrictionofto{\catindexof{\nodes}}{\secedge\cap\secnodes}
            ,\catvariableof{\secedge/\secnodes}}} \neq 0
    } \, .
\end{align*}
One can understand the guess of an assignment $\catindexof{\node}$ to a variable $\catvariableof{\node}$, as it is done during backtracking search, as an inclusion of a constraint
\begin{align*}
    \kcoreofat{\{\node\}}{\catvariableof{\node}}
    = \onehotmapofat{\catindexof{\node}}{\catvariableof{\node}} \, .
\end{align*}
Therefore, constraint propagation \algoref{alg:knowledgePropagation} can be integrated with backtracking search, with iterations between propagations of knowledge and guessing of additional variables.
%The message cores $\kcoreof{\edge}$ are subset encoding of possible local choices, according to which variables can be assigned.

\subsection{Discussion}

In our perspective, the number of possible sub-networks to be contracted is the bottleneck, instead of the demand to perform such contractions.
%The number of sub-networks grows as binomial coefficients.
The choice and the scheduling of these sub-contractions is most important for the design of efficient inference schemes.
